% GENERATED FILE. Edit canonical lessons, then run npm run generate:books.
% canonical-source-hash: fnv1a64:bfacac41e149f323
% canonical-lessons: PA-C359-sabak, PA-C359-abhias, PA-C359-mihnat, PA-C359-dhian, PA-C359-sudhaarna

\chapter{Studying Hard}
\label{ch:pa-a2-studying-hard}
\hlchaptermodality{359}

\begin{chapteropening}
\textbf{By the end of this chapter:} \emph{I can say how I study and practise.}

Complete the last lesson of chapter 359: I can say how I study and practise.
\end{chapteropening}

\section[\texorpdfstring{\pa{ਸਬਕ} (sabak)}{ਸਬਕ (sabak)}]{\pa{ਸਬਕ} (sabak) — a lesson to learn}

\label{lesson:PA-C359-sabak}



\begin{quote}
Before the new one: say the Punjabi for history, then the Punjabi for to teach a subject.
\end{quote}

\subsection*{You'll want to know: \pa{ਸਬਕ}}
\textbf{\pa{ਸਬਕ}} — \emph{sabak} — \textquotedblleft{}a lesson to learn\textquotedblright{}.

\begin{quote}
\pa{ਅੱਜ} \pa{ਦਾ} \pa{ਸਬਕ} \pa{ਛੋਟਾ} \pa{ਹੈ।}
\end{quote}

\begin{quote}
\emph{ajj dā sabak chhoṭā hai.}
\end{quote}

\begin{quote}
Today's lesson is short.
\end{quote}

A masculine noun.

\begin{cousinweb}
\textbf{\pa{ਸਬਕ}} (\emph{sabak}) is a loan, usually traced through Persian to Arabic \emph{sabaq}, \textquotedblleft{}lesson\textquotedblright{}.
\end{cousinweb}

\begin{grammarlens}[title={What you've built}]
One more word for this chapter.
\end{grammarlens}

\subsection*{Your turn}
\begin{itemize}
\raggedright
  \item \emph{Say it:} \emph{sabak}
  \item \emph{Say it:} \emph{sabak}, once more
  \item \emph{Say it:} \emph{ajj dā sabak chhoṭā hai.} — Today's lesson is short.
  \item \emph{Recall:} say the Punjabi for history, then the Punjabi for to teach a subject, then say \emph{sabak} again
\end{itemize}

\begin{tcolorbox}[breakable,colback=teal!4,colframe=teal!35!black,title={Before you move on}]
What does \pa{ਸਬਕ} mean? (\textquotedblleft{}A lesson to learn\textquotedblright{}.) Say it once more.
\end{tcolorbox}
\section[\texorpdfstring{\pa{ਅਭਿਆਸ} (abhiās)}{ਅਭਿਆਸ (abhiās)}]{\pa{ਅਭਿਆਸ} (abhiās) — practice; drill}

\label{lesson:PA-C359-abhias}



\begin{quote}
Before the new one: say the Punjabi for to teach a subject, then the Punjabi for a lesson to learn.
\end{quote}

\subsection*{You'll want to know: \pa{ਅਭਿਆਸ}}
\textbf{\pa{ਅਭਿਆਸ}} — \emph{abhiās} — \textquotedblleft{}practice; drill\textquotedblright{}.

\begin{quote}
\pa{ਮੈਂ} \pa{ਹਰ} \pa{ਰੋਜ਼} \pa{ਅਭਿਆਸ} \pa{ਕਰਦਾ} \pa{ਹਾਂ।}
\end{quote}

\begin{quote}
\emph{maiṁ har roz abhiās kardā hāṁ.}
\end{quote}

\begin{quote}
I practise every day.
\end{quote}

A masculine noun.

\begin{cousinweb}
\textbf{\pa{ਅਭਿਆਸ}} (\emph{abhiās}) is taken over from Sanskrit \emph{abhyāsa}, \textquotedblleft{}repetition\textquotedblright{}, a tatsama (Monier-Williams).
\end{cousinweb}

\begin{grammarlens}[title={What you've built}]
One more word for this chapter.
\end{grammarlens}

\subsection*{Your turn}
\begin{itemize}
\raggedright
  \item \emph{Say it:} \emph{abhiās}
  \item \emph{Say it:} \emph{abhiās}, once more
  \item \emph{Say it:} \emph{maiṁ har roz abhiās kardā hāṁ.} — I practise every day.
  \item \emph{Recall:} say the Punjabi for to teach a subject, then the Punjabi for a lesson to learn, then say \emph{abhiās} again
  \item \emph{Write it:} copy \textbf{\pa{ਅਭਿਆਸ}} once, checking the whole word
\end{itemize}

\begin{tcolorbox}[breakable,colback=teal!4,colframe=teal!35!black,title={Before you move on}]
What does \pa{ਅਭਿਆਸ} mean? (\textquotedblleft{}Practice; drill\textquotedblright{}.) Say it once more.
\end{tcolorbox}
\section[\texorpdfstring{\pa{ਮਿਹਨਤ} (mihnat)}{ਮਿਹਨਤ (mihnat)}]{\pa{ਮਿਹਨਤ} (mihnat) — hard work; effort}

\label{lesson:PA-C359-mihnat}



\begin{quote}
Before the new one: say the Punjabi for a lesson to learn, then the Punjabi for practice.
\end{quote}

\subsection*{You'll want to know: \pa{ਮਿਹਨਤ}}
\textbf{\pa{ਮਿਹਨਤ}} — \emph{mihnat} — \textquotedblleft{}hard work; effort\textquotedblright{}.

\begin{quote}
\pa{ਉਹ} \pa{ਬਹੁਤ} \pa{ਮਿਹਨਤ} \pa{ਕਰਦੀ} \pa{ਹੈ।}
\end{quote}

\begin{quote}
\emph{oh bahut mihnat kardī hai.}
\end{quote}

\begin{quote}
She works very hard.
\end{quote}

A feminine noun.

\begin{cousinweb}
\textbf{\pa{ਮਿਹਨਤ}} (\emph{mihnat}) is a loan, usually traced through Persian to Arabic \emph{miḥna}, \textquotedblleft{}trial\textquotedblright{}.
\end{cousinweb}

\begin{grammarlens}[title={What you've built}]
One more word for this chapter.
\end{grammarlens}

\subsection*{Your turn}
\begin{itemize}
\raggedright
  \item \emph{Say it:} \emph{mihnat}
  \item \emph{Say it:} \emph{mihnat}, once more
  \item \emph{Say it:} \emph{oh bahut mihnat kardī hai.} — She works very hard.
  \item \emph{Recall:} say the Punjabi for a lesson to learn, then the Punjabi for practice, then say \emph{mihnat} again
  \item \emph{Write it:} copy \textbf{\pa{ਮਿਹਨਤ}} once, checking the whole word
\end{itemize}

\begin{tcolorbox}[breakable,colback=teal!4,colframe=teal!35!black,title={Before you move on}]
What does \pa{ਮਿਹਨਤ} mean? (\textquotedblleft{}Hard work; effort\textquotedblright{}.) Say it once more.
\end{tcolorbox}
\section[\texorpdfstring{\pa{ਧਿਆਨ} (dhiān)}{ਧਿਆਨ (dhiān)}]{\pa{ਧਿਆਨ} (dhiān) — attention}

\label{lesson:PA-C359-dhian}



\begin{quote}
Before the new one: say the Punjabi for practice, then the Punjabi for hard work.
\end{quote}

\subsection*{You'll want to know: \pa{ਧਿਆਨ}}
\textbf{\pa{ਧਿਆਨ}} — \emph{dhiān} — \textquotedblleft{}attention\textquotedblright{}.

\begin{quote}
\pa{ਜਮਾਤ} \pa{ਵਿੱਚ} \pa{ਧਿਆਨ} \pa{ਨਾਲ} \pa{ਸੁਣੋ।}
\end{quote}

\begin{quote}
\emph{jamāt vich dhiān nāl suṇo.}
\end{quote}

\begin{quote}
Listen carefully in class.
\end{quote}

A masculine noun.

\begin{cousinweb}
\textbf{\pa{ਧਿਆਨ}} (\emph{dhiān}) is taken over from Sanskrit \emph{dhyāna}, \textquotedblleft{}meditation, thought, reflection\textquotedblright{}, a tatsama (Monier-Williams); the everyday sense \textquotedblleft{}attention\textquotedblright{} is not among them.
\end{cousinweb}

\begin{grammarlens}[title={What you've built}]
One more word for this chapter.
\end{grammarlens}

\subsection*{Your turn}
\begin{itemize}
\raggedright
  \item \emph{Say it:} \emph{dhiān}
  \item \emph{Say it:} \emph{dhiān}, once more
  \item \emph{Say it:} \emph{jamāt vich dhiān nāl suṇo.} — Listen carefully in class.
  \item \emph{Recall:} say the Punjabi for practice, then the Punjabi for hard work, then say \emph{dhiān} again
  \item \emph{Write it:} copy \textbf{\pa{ਧਿਆਨ}} once, checking the whole word
\end{itemize}

\begin{tcolorbox}[breakable,colback=teal!4,colframe=teal!35!black,title={Before you move on}]
What does \pa{ਧਿਆਨ} mean? (\textquotedblleft{}Attention\textquotedblright{}.) Say it once more.
\end{tcolorbox}
\section[\texorpdfstring{\pa{ਸੁਧਾਰਨਾ} (sudhārnā)}{ਸੁਧਾਰਨਾ (sudhārnā)}]{\pa{ਸੁਧਾਰਨਾ} (sudhārnā) — to correct; to put right}

\label{lesson:PA-C359-sudhaarna}



\begin{quote}
Before the new one: say the Punjabi for hard work, then the Punjabi for attention.
\end{quote}

\subsection*{You'll want to know: \pa{ਸੁਧਾਰਨਾ}}
\textbf{\pa{ਸੁਧਾਰਨਾ}} — \emph{sudhārnā} — \textquotedblleft{}to correct; to put right\textquotedblright{}.

\begin{quote}
\pa{ਅਧਿਆਪਕ} \pa{ਮੇਰੇ} \pa{ਗ਼ਲਤ} \pa{ਸ਼ਬਦ} \pa{ਸੁਧਾਰਦਾ} \pa{ਹੈ।}
\end{quote}

\begin{quote}
\emph{adhiāpak mere ghalat shabad sudhārdā hai.}
\end{quote}

\begin{quote}
The teacher corrects my wrong words.
\end{quote}

Past (I …): \pa{ਮੈਂ} \pa{ਗ਼ਲਤ} \pa{ਸ਼ਬਦ} \pa{ਸੁਧਾਰੇ}, \emph{maiṁ ghalat shabad sudhāre}.

\begin{cousinweb}
\textbf{\pa{ਸੁਧਾਰਨਾ}} (\emph{sudhārnā}) is inherited from a reconstructed \emph{*śuddhakāra}, \textquotedblleft{}making clean\textquotedblright{}; CDIAL lists Punjabi \emph{sudharnā}, \textquotedblleft{}to repair\textquotedblright{}, under it (CDIAL 12521). CDIAL does not give the sense \textquotedblleft{}to correct\textquotedblright{}.
\end{cousinweb}

\begin{grammarlens}[title={What you've built}]
One more word for this chapter.
\end{grammarlens}

\subsection*{Your turn}
\begin{itemize}
\raggedright
  \item \emph{Say it:} \emph{sudhārnā}
  \item \emph{Say it:} \emph{sudhārnā}, once more
  \item \emph{Say it:} \emph{adhiāpak mere ghalat shabad sudhārdā hai.} — The teacher corrects my wrong words.
  \item \emph{Recall:} say the Punjabi for hard work, then the Punjabi for attention, then say \emph{sudhārnā} again
  \item \emph{Write it:} copy \textbf{\pa{ਸੁਧਾਰਨਾ}} once, checking the whole word
\end{itemize}

\begin{tcolorbox}[breakable,colback=teal!4,colframe=teal!35!black,title={Before you move on}]
What does \pa{ਸੁਧਾਰਨਾ} mean? (\textquotedblleft{}To correct; to put right\textquotedblright{}.) Say it once more.
\end{tcolorbox}
